% Section 6: Discussion
% Paper: "Just a Glimpse: Temporal Wearing as a Design Dimension for Transient Mixed Reality"

\section{Discussion}

Our findings across both study phases suggest that temporal wearing is not merely a shorter version of session wear but a qualitatively different mode of interaction with its own behavioral patterns, social dynamics, and design requirements. In this section, we reflect on wearing duration as a design dimension, discuss what the baseline comparison revealed, articulate six design principles derived from our empirical findings, consider applications beyond mixed reality, and acknowledge limitations.

\subsection{Wearing Duration as a Design Dimension}

The dominant trajectory of MR development treats progress as longer, more immersive, more seamless wearing. Each generation of HMDs is lighter, more comfortable, and more capable---all in service of extending wearing duration. Apple Vision Pro's ``all-day'' comfort, Meta Quest 3's extended battery life, and the industry's pursuit of lightweight AR glasses all reflect this logic: the ideal device is one you never take off. Our findings suggest a complementary perspective---not that sustained wear is wrong, but that wearing duration is a consequential design dimension that has been undertheorized, and that the brief end of the duration spectrum has distinct properties worth designing for.

The lab study (Phase~1) provides initial evidence for this claim. When participants used the same content through temporal wearing versus session wear, the two modes appeared to produce different interaction patterns: temporal wearing conditions elicited more frequent re-engagement cycles, and participants reported [PLACEHOLDER: higher/comparable] social comfort with temporal wearing forms. The deployment (Phase~2) extends this finding into a naturalistic setting, where the emergent behaviors---double-takes, hand-offs, narration, orbiting---suggest that temporal wearing invites interaction patterns that session wear does not readily afford.

When participants performed double-takes, hand-offs, narrations, and orbits, they were not ``failing to sustain engagement''---they were engaging in precisely the way temporal wearing invites. The efficiency ratio data (Section~5.1) indicates that temporal wearing can deliver augmented experiences with minimal mechanical overhead when transition costs are minimized: with sub-second onset times, even a three-second glimpse devotes the majority of the encounter to augmented content. This is not a compromise but a design choice---and one that appears to bring benefits in social comfort, shareability, and naturalness that sustained wear may not match.

Our findings also suggest that MR experience may not accumulate only linearly. The orbit behavior (Section~5.2) indicates that temporal wearing can support a \textit{cumulative} model of experience, where meaning builds across multiple brief encounters rather than within a single extended session. This pattern resonates with how people experience many cultural artifacts---returning to a painting, re-reading a poem, revisiting a scenic view---and suggests that MR content designed for temporal wearing may prove engaging in ways that differ from content designed for sustained immersion.

\subsection{What the Baseline Comparison Revealed}

The session-wear baseline in Phase~1 serves a critical function: it anchors our claims about temporal wearing in a direct comparison rather than relying solely on self-reported contrasts with prior HMD experience. [PLACEHOLDER: Discuss key findings from the baseline comparison. Expected topics:]

The comparison suggests several patterns worth noting. First, transition cost: the session-wear condition required [PLACEHOLDER: X] seconds for donning (strap adjustment, IPD calibration, orientation) compared to [PLACEHOLDER: <2] seconds for all temporal wearing conditions. This [PLACEHOLDER: order-of-magnitude] difference appears to be not just quantitative but qualitative---participants described the session-wear don as a ``setup'' process and the temporal wearing don as a ``gesture.''

Second, social comfort: participants reported [PLACEHOLDER: lower] social comfort in the session-wear condition, even in the controlled lab setting. Bystander companions reported [PLACEHOLDER: findings about watching someone in an HMD vs. watching someone peek through a handheld viewer]. These findings are consistent with the social acceptability literature~\cite{koelle2020social} but extend it by providing a direct within-subjects comparison between temporal and session wearing of the \textit{same content}.

Third, cognitive load: the NASA-TLX comparison revealed [PLACEHOLDER: findings]. We note that [PLACEHOLDER: any surprising findings, e.g., experience quality ratings were comparable despite the much briefer exposure in temporal wearing conditions, suggesting that temporal wearing's information density per second may compensate for shorter total exposure].

These baseline findings should be interpreted with caution. The lab setting is artificial, the session-wear condition used a specific HMD model, and the content was designed for temporal wearing (brief, self-contained) rather than for sustained session use. A fairer comparison might use content optimized for each wearing mode. Our findings indicate that temporal wearing differs from session wear on measurable dimensions, but the relative advantages of each mode likely depend on context, content, and user goals.

\subsection{Design Principles for Temporal Wearing}

Drawing on our design process and empirical findings from both study phases, we propose six design principles for temporal wearing devices. Each principle is grounded in specific findings; we note the supporting evidence and moderate our claims accordingly.

\subsubsection{P1: Design for Discovery.}
The device should invite curiosity before any interaction begins. Our deployment data suggests that the \textit{attract} phase of the glimpse cycle (Section~3.5) plays an important role: visitors who noticed the device---through its physical presence, through signage, or through observing another visitor's reaction---appeared more likely to engage than those who passed by unaware. The honeypot effect~\cite{wouters2016honeypot}---where one person's use attracts others---emerged as a discovery mechanism in [PLACEHOLDER: \%] of first interactions (Phase~2 video coding).

Design implications: Temporal wearing devices should be physically present and visually legible. Their form should communicate ``look through me'' through affordances~\cite{norman2013design}: lens-like apertures, viewfinder shapes, or transparent optics that hint at something beyond.

\subsubsection{P2: Instant-On.}
The transition from not-wearing to augmented perception should take less than two seconds. Our Phase~1 data indicates that transition cost is a primary mechanical correlate of re-engagement frequency: participants re-engaged more frequently with prototypes that offered near-instantaneous onset. Phase~2 deployment data is consistent with this pattern ($\rho$ = [PLACEHOLDER], $p$ = [PLACEHOLDER]).

Design implications: Displays should be always-on or wake from sleep in under 100ms. Tracking should be pre-initialized or use inside-out methods that do not require calibration. A useful heuristic: if the user can complete the don gesture faster than the device can deliver content, the device is too slow.

\subsubsection{P3: Graceful Doff.}
Exiting the augmented experience should feel natural, not abrupt or punishing. Interview participants in both phases contrasted temporal wearing with VR experiences where ``taking off the headset'' felt like ``being ejected from'' the experience. Temporal wearing doffing should feel like lowering binoculars---a natural return to unmediated perception.

Design implications: Avoid content that requires the user to complete an action before doffing. Design content that is complete at any moment of exit---every frame should be a valid stopping point. Optical see-through displays inherently support this by allowing the real world to persist behind the augmented layer.

\subsubsection{P4: Social Legibility.}
Bystanders should be able to perceive what the user is doing, why, and that it will be brief. Our social comfort data from both phases suggests that temporal wearing's social advantage depends on legibility: participants felt comfortable using the devices when the interaction ``looked normal'' and was visibly brief. The Phase~1 baseline comparison indicates that session-wear's lower social comfort scores may stem partly from reduced legibility---bystanders found it harder to interpret what the HMD user was doing and when they would finish.

Design implications: The donning gesture should reference familiar actions: raising binoculars, looking through a viewfinder, leaning toward a telescope. The device's form should communicate its purpose even to people who have never seen it before.

\subsubsection{P5: Invitation to Return.}
Content should reward re-engagement rather than penalize brevity. The orbit behavior (Phase~2) suggests that temporal wearing naturally invites return visits---but our interview theme of ``wanting more (in the right way)'' (Section~5.2) captures a critical distinction: temporal wearing should create anticipatory desire (``I want to come back'') rather than deprivation (``I wish it lasted longer'').

Design implications: Content should vary across encounters---through spatial variation, temporal variation, progressive revelation, or stochastic elements. The ideal is a ``generous preview'': each glimpse is complete and satisfying in itself, yet hints at more to be discovered.

\subsubsection{P6: Shared Glimpsing.}
Temporal wearing devices should enable, not obstruct, social interaction around the moment of augmented perception. The narrator, hand-off, and social permission behaviors (Phase~2) suggest that temporal wearing can produce shared MR experiences without requiring multiple synchronized devices. This shareability appears to be one of temporal wearing's most distinctive properties compared to session wear.

Design implications: Design for passability: handheld devices should be easy to pass between people. Design for narration: the user should be able to speak while wearing the device and should maintain awareness of companions. Design for spectatorship: bystanders should find it interesting to watch someone use the device.

\subsection{Temporal Wearing Beyond Mixed Reality}

While we have developed temporal wearing in the context of MR, the design dimension may extend to other domains of wearable and ubiquitous computing where brief, bounded device encounters offer advantages over sustained use.

\paragraph{Internet of Things.}
Many interactions with smart home devices are temporal: a brief check of a thermostat display, a momentary glance at an energy monitor. Designing these interactions explicitly with temporal wearing principles---instant-on, graceful exit, invitation to return---could improve their integration into domestic routines.

\paragraph{Health Monitoring.}
For many health-conscious individuals, a brief daily check---temporal wearing of a diagnostic device---may be more appropriate than continuous monitoring. Designing health devices for temporal wearing could reduce the burden and stigma of health monitoring while still providing actionable data.

\paragraph{Accessibility.}
Some users of assistive technology prefer intermittent use---engaging with assistive technology only in specific contexts~\cite{profita2016acceptability}. Temporal wearing offers a design framework for assistive devices that support this pattern: rapid onset when needed, graceful doff when not, and social legibility that reduces stigma.

\subsection{Limitations and Future Work}

Our work has several limitations that suggest directions for future research.

\paragraph{Lab Study Limitations.}
The Phase~1 lab study (N=20) provides controlled comparisons but in an artificial setting. The session-wear baseline used a single HMD model with content designed for temporal wearing, which may disadvantage the session-wear condition. Future work should compare temporal and session wearing with content optimized for each mode. The sample size, while adequate for within-subjects comparisons, limits generalizability; larger studies with demographic stratification would strengthen these findings.

\paragraph{Single Deployment Context.}
Our Phase~2 findings are drawn from a single deployment at [VENUE PLACEHOLDER]. While we selected this venue for its representative qualities, temporal wearing may manifest differently in other contexts---outdoor environments, retail settings, educational institutions, or domestic spaces. Future work should examine temporal wearing across diverse deployment contexts.

\paragraph{Content Effects.}
All prototypes presented the same augmented content, allowing us to isolate form factor effects. However, content type likely interacts with temporal wearing behavior: narrative content may encourage longer glimpses, while abstract visual content may support faster cycling. Systematically varying content alongside form factor would clarify these interactions.

\paragraph{Novelty Effect.}
Both study phases captured participants' first encounters with temporal wearing prototypes. The novelty of the devices may have inflated engagement metrics and positive responses. Longitudinal deployments---measuring how temporal wearing behaviors evolve over repeated visits across days or weeks---would address whether the patterns we observed are durable. We note, however, that the historical persistence of temporal wearing devices (coin-operated viewfinders, View-Masters) suggests that the practice itself may be robust to novelty decay, even if specific devices lose appeal.

\paragraph{Head-Mounted Temporal Wearing.}
Our prototype set did not include a head-mounted temporal wearing device. While head-mounted devices inherently incur higher transition costs, it is possible to design lightweight masks or flip-up visors that support temporal wearing from a head-mounted position. Future work should explore whether head-mounted form factors can achieve the sub-two-second transition costs that our findings suggest are important for temporal wearing.

\paragraph{Longitudinal Content Design.}
We proposed a layered revelation content model but did not empirically test it. Designing and deploying content specifically structured for cumulative glimpses, and measuring how visitor engagement evolves across multiple return visits, is a natural next step.
